\section{Version history}
\label{app:versions}

\paragraph{\manuscriptversion\ (\manuscriptdate).}
This version follows a complete mathematical, numerical and formal review of the analysis. The main changes from Version 1 are:
\begin{itemize}[leftmargin=*, itemsep=1pt]
  \item The Version 1 tightening magnitudes (about $0.109$\,eV with SPT and DESI, and about $0.012$\,eV with the Planck-subset baseline) were derived from single, unconverged chains. They are withdrawn as unsupported. They have not been shown to be false. No replacement magnitude is given.
  \item Localization previously inverted each block of the covariance separately. It now uses the exact signed allocation of \cref{sec:ledger}. The audits were re-run, including cosmological audits with a shared Boltzmann calculation and multi-start fits. The claim that high-$\ell$ TT dominates the mismatch in general is replaced by the baseline- and direction-dependent statement of \cref{sec:results-audits}.
  \item The fixed-spectrum profiled audit now releases the shared parameters in the reference fit and uses multiple starting points. Its values change from $859.6$/$369.5$ to $849.6$/$414.4$.
  \item The boundary-mode discussion now states its hypotheses explicitly. It shows that a boundary mode requires a nonpositive combined mean, which two constraints with positive means cannot produce, corrects the claim that tightening generally increases boundary-mode frequency (\cref{prop:sweep}), and adds the finite-upper-gate and readout results. The Lean library grew from two lemmas to 97 audited declarations.
  \item Chain summaries now use the inverse weighted empirical CDF with exact holding-time weights, together with chronological split diagnostics. For the same stored samples, this changes the archived SPT D1 + DESI 95th percentile from $0.085$ to $0.098$\,eV, and it shows that none of the archived chains is converged (\cref{tab:archived-chains}).
  \item The abstract and main text were rewritten. The audit figures and tables are now generated directly from the verified run bundles, and the two-constraint figure is re-evaluated from the analytic formulas.
\end{itemize}

\paragraph{Version 1 (2 March 2026).}
First public version, posted on Research Square \cite{tsiokos2026_neutrino_v1}.
